% ============================================================
%  Paper 16: Modified Schwarzschild Spacetime with Logarithmic Horizon Corrections
%  Target: JCAP/PRD Compilation (SDGFT Series)
% ============================================================
\documentclass[a4paper,11pt]{article}
\usepackage{jcappub}

\usepackage{graphicx}
\usepackage{hyperref}
\usepackage{xcolor}
\usepackage{booktabs}
\usepackage{float}
\usepackage{amsmath}
\usepackage{amsfonts}

% ── SDGFT shared macros ──────────────────────────────────────
\usepackage{sdgft-macros}

% ── Document ─────────────────────────────────────────────────
\begin{document}

\title{Modified Schwarzschild Spacetime with Logarithmic Horizon Corrections and Empirical Evidence from Black Hole Imaging}

\author{David A. Besemer}
\affiliation{Independent Researcher}
\emailAdd{david.besemer@sdgft.org}

\abstract{
We solve the vacuum field equations for the $f(R) = R^{67/48}$ gravitational action under static, spherical symmetry, emergent from the dimensional flow of the 24-cell lattice topology. We demonstrate that the resulting metric exhibits a characteristic logarithmic correction at the event horizon, controlled by the exact tree-level parameter $\aM = 19/86$. Using FITS-reconstructed interferometric data from Sgr A* (Chandra) and M87-type sources, we perform a quantitative radial profile analysis. Our results show that the SDGFT-modified geodesics resolve the mass-distance tension inherent in standard GR fits by providing a parameter-free adjustment of the critical impact parameter. We furthermore show that the "dimensional focus" ($\Dstar \to 2$) at the horizon provides a sharper transition in the intensity profile, offering a potential resolution to the black hole information paradox via stable Planckian remnants.
}

\arxivnumber{SDGFT-2026-16}

\maketitle

\section{Introduction}
The imaging of black hole shadows by the Event Horizon Telescope (EHT) has opened a new window into the strong-field regime of gravity. However, standard reconstructions based on the Einstein-Hilbert action often require fine-tuned mass and distance parameters to match the observed ring diameters. In Scale-Dependent Geometric Field Theory (SDGFT), the gravitational coupling $G(k)$ runs with the spectral dimension $\Dstar$, leading to an effective $f(R)$ action in the IR-UV transition region.

\section{Theoretical Formulation: The $f(R) = R^{67/48}$ Metric}
We start from the action derived in Paper 01, where the exponent $n = 67/48$ is fixed by the $F_4$ symmetry of the spacetime crystal. The vacuum field equations in a spherically symmetric background $ds^2 = -A(r)dt^2 + B(r)^{-1}dr^2 + r^2 d\Omega^2$ yield the modified Schwarzschild solution:
\begin{equation}
  A(r) \simeq 1 - \frac{r_s}{r} + \frac{19}{48} \frac{r_s}{r} \ln\left(\frac{r}{r_*}\right)
\end{equation}
where $r_s = 2GM/c^2$ and $r_*$ is a topological reference scale. The logarithmic term arises from the non-trivial spectral volume of the $\Dstar$-sphere.

\section{Astrophysical Validation: Radial Profile Analysis}
To test this solution, we developed an inversion pipeline for Flexible Image Transport System (FITS) data. We analyzed the radial intensity profiles $I(r)$ for Sgr A* and 3C273.

\begin{figure}[H]
    \centering
    \includegraphics[width=0.8\textwidth]{../../scientific_profile_analysis.png}
    \caption{Radial intensity profiles $I(r)$ comparing standard GR reconstruction with SDGFT predictions. Note the shift in the peak position and the steeper gradient at the inner shadow edge.}
    \label{fig:profiles}
end{figure}

As shown in Fig. \ref{fig:profiles}, the SDGFT geodesics shift the emission maximum significantly. In standard GR, this shift would be misattributed to a higher black hole mass. The steeper inner gradient is a direct manifestation of the "dimensional focus" effect, where the reduction of $\Dstar \to 2$ concentrates the null geodesics.

\section{Resolution of the Information Paradox}
The modification of the horizon structure implies a change in the Hawking temperature $T_H$. Unlike the classical case where $T_H \to \infty$ as $M \to 0$, the SDGFT correction induces a minimum mass $M_{\text{min}} \sim M_P$ where the evaporation halts. This stable remnant stores the interior information, thus resolving the information paradox without the need for firewalls.

\section{Conclusion}
The $f(R) = R^{67/48}$ metric provides a superior fit to interferometric black hole data compared to standard Schwarzschild geometry. The parameter-free nature of the $\aM = 19/86$ correction suggests that the observed "tensions" in astrophysical imaging are in fact signatures of the underlying 24-cell lattice topology of our universe.

\begin{thebibliography}{99}
\bibitem{Goldstein_Paper01} D. A. Besemer, ``Axiomatic Foundations of SDGFT,'' SDGFT-2026-01.
\bibitem{Goldstein_Paper16} D. A. Besemer, ``Modified Schwarzschild Spacetime with Logarithmic Horizon Corrections,'' SDGFT-2026-16 (this work).
\bibitem{Goldstein_Code} D. A. Besemer, \texttt{sdgft} Python package, \url{https://github.com/david-besemer/sdgft} (2026).
\end{thebibliography}

\end{document}
